\chapter{What the complete project establishes}
\section{There is no single controller ranking across these studies}
The project does not support the statement that base PI is best for everything. It is especially strong on the frozen-controller, five-second no-load acceleration test. At 9250 rpm with load pulses, retuned Poly2 improves both RMSE and fuel ringing over base PI. The original PILCO-style policy improves RMSE further but produces much more ringing. Smooth PILCO substantially reduces that ringing at the cost of tracking performance. In the latest combined cycle, DOB98 with retuned PI gives the lowest full-cycle RMSE among the revised four-controller comparison.

These conclusions coexist because the tasks and objectives differ. Ramp tracking requires coordinated acceleration across operating points. Load rejection requires restoring a fixed operating point after an unknown resisting torque. A fuel-reversal penalty discourages oscillatory correction but does not directly minimize peak speed error, peak fuel slew, or recovery time. A report that declares one universal winner by averaging incompatible scenarios would conceal the central result of the work.

\section{Why added sophistication sometimes failed to help}
The first reason is information. Static nominal feedforward cannot cancel an unknown load at an unchanged reference. At fixed conditions, its constant contribution can be absorbed by an appropriately initialized integral state. The second reason is structure: measured-state feedforward changes the feedback loop, while inverse-coordinate PI introduces acceleration-sensitive action. An accurate regression is not sufficient if these changes are poorly scaled.

The third reason is objective mismatch. Original PILCO learned gains under a local load-pulse objective that rewarded tracking and fuel magnitude. It did not penalize repeated reversals until the ringing term was added. Neither policy was trained to optimize the 9000--10000 rpm ramp. The fourth reason is distribution shift: a map or dynamics model fitted near a limited speed/ambient region cannot be presumed valid over every later trajectory. The broad ambient cycle made this issue visible through range violations and fallback.

The fifth reason is baseline strength and tuning attribution. The inherited PI already closes the loop around the nonlinear plant. R3's target result is nearly reproduced by a matched filtered PID, indicating that architecture and tuning account for much of the benefit at that point. The latest DOB comparison changes both observer structure and PI gains, but preserves an intermediate DOB-plus-original-PI result. That intermediate result is essential for separating observer benefit from subsequent retuning.

\section{The sensor audit resolves one hypothesis}
The explicit same-sensor rerun reproduces the preceding speed, sensed-speed, and fuel trajectories. Unequal sensor dynamics therefore do not explain the earlier ranking. The common sensor still matters because the controllers respond differently to its lag. At 200 rpm/s, keeping sensed speed close to the target can place actual speed above it. The base controller's feedback lag can partly offset that lead in the actual-speed metric.

This is an interpretation consistent with the recorded signal differences and a first-order lag calculation. It is not a proof that sensor lag is beneficial, nor a reason to add delay deliberately. A zero-lag or delay-compensated controlled ablation would be a new experiment and was not performed for this report. The appropriate conclusion is to define the scored signal clearly and preserve common sensing when comparing controllers.

\section{A hierarchy of validation claims}
Algebraic replay verifies that the implementation follows its stated recurrence. Frequency-response identities verify the constructed filtered inverse. Monte Carlo comparisons in the archived PILCO studies verify moment calculations on the same GP model. None alone validates the nonlinear turbine model. Saved nonlinear simulation traces test closed-loop behavior for their specified scenarios, subject to solver convergence and map limitations. Frozen-parameter amplitude checks supply limited transfer evidence. Hardware testing would address a different and currently absent level of validation.

The latest 9800-rpm model's excellent free-run fit is local: the identification speeds vary only a few rpm. The combined cycle covers a much larger range and briefly reaches fuel saturation. Its observed success is valuable empirical evidence, but the local fit cannot be promoted into a global error bound. Similarly, positive surge margin in a converged nominal simulation is not a certified engine-operating envelope.

\section{What remains untested}
The archive does not establish robustness to sensor noise, bias, drift, communication delay, actuator-rate limits, uncertain component maps, deterioration, or different shaft inertia. It does not optimize all controller families with an equal tuning budget on a common training distribution. It does not establish global stability under fallback switching, gain scheduling, or saturation. It does not report real-time execution benchmarks or hardware fuel savings.

A future program could first declare a common multi-condition training/validation split, then retune competing architectures under identical constraints. Noise and actuator limits could be introduced separately, with failure criteria fixed before examining results. A local-model schedule could be compared with the single 9800-rpm observer. Such proposals are future work, not simulations added to this retrospective report.

\section{Lessons for an undergraduate control engineer}
First, verify the experiment before interpreting the controller. The failed chirp and PRBS are methodological evidence, not inconveniences to discard. Second, specify what information a predictor receives: one-step and free-running accuracy answer different questions. Third, inspect inversion properties rather than treating high fit as a license to invert. Fourth, write the complete control law including sensor, initialization, saturation, and guards. Fifth, report multiple performance dimensions and their windows.

The most useful achievement of this project is therefore a connected set of increasingly discriminating experiments. Each disappointing result narrows the next hypothesis. The progression from ARX equilibrium correction to local DOB98, and from inverse-GP accuracy to polynomial and PID challengers, illustrates how practical control research improves through explicit assumptions, matched evidence, and honest limits.

\chapter{Reproducibility, presentations, and preservation}
\section{Three different reproduction tasks}
Document reproduction means compiling the supplied LaTeX source and included figures. It requires no MATLAB and no new simulation. Analysis reproduction means recomputing metrics or figures from saved trajectories, preserving the original acceptance masks and scoring windows. Experiment reproduction means rerunning plant models or training scripts with the original dependencies. These tasks are deliberately separated so that reading or rebuilding the report cannot silently launch a new experiment.

The standalone source package contains the assembled manuscript, imported vector figures, generated tables, explanatory chapters, archived narrative evidence, and a SHA-256 evidence index. Source paths in the index are relative to the original workspace. Raw MAT files and every full trajectory need not be duplicated in a document bundle; the study directories and repository dataset catalogs identify their locations. The evidence index records exactly which files supported this edition.

\section{Animations are views of recorded data}
The project produced progressively improved controller-comparison videos. Later versions use distinct windows for each controller, synchronized time, visible external load, common axes where appropriate, and speed-error/fuel plots that reveal differences hidden by the 9000--10000 rpm scale. The rotor depiction is a visual aid with scaled apparent rotation, not a blade-resolved simulation. Video interpolation or playback speed must not be confused with the 15 ms simulation sample interval.

The historical PILCO, smoother-policy, no-load ramp, equal-sensor, combined-cycle, and DOB98 videos document their corresponding archived comparisons. Removing previous PILCO from the latest four-controller video does not erase its earlier study; its results remain necessary to explain why a ringing penalty was introduced. The report preserves this distinction between the current demonstration lineup and the full research history.

\section{Repository and report lineage}
The earlier observer studies were distributed through the Research repository. The separate TMATS-GPR-Poly2-Control repository originally covered the GPR-to-Poly2 phase through September 15. This report crosses that split and includes the October studies. A repository update should therefore link the complete thesis, preserve historical reports, include the latest controller source and results, and explain which files are data, which are simulations, and which merely build presentation artifacts.

The revised September paper corrected signal timing, anti-windup definitions, fallback semantics, and interpretation of actuator demand after AI referee-role reviews. Those corrections are retained here. The October ramp report's compact policy notation used an ambiguous speed symbol; this edition explicitly uses $N_s$ for feedback and $N$ for actual-speed scoring. The distinction changes the explanation, not the archived trajectory. Historical file names and dates are preserved even where the narrative has been clarified.

\section{How to extend the record responsibly}
Keep earlier accepted and failed runs rather than overwriting them. Save the actual applied reference, load, fuel, sensed and mechanical speed, validity flags, solver diagnostics, and controller internal states. Record gain selection separately from held-out evaluation. Hash the configuration and include units. When a metric definition changes, recompute matched comparators or keep the resulting table separate. These practices make future results comparable to the present record without implying that every older experiment used the same protocol.
